% !TeX root = ../document_maitre.tex

Ces innovations technologiques promettent des évolutions importantes pour les \glspl{sgc} et le monde professionnel concerné. Elles ont pour but d'accompagner les professionnels face aux nouveaux défis et enjeux de la gestion des collections. Pour cela, on a identifié trois axes d'améliorations :

L'axe de l'accessibilité aux données stockées, par une recherche plus intuitive et puissante qu'est la \guillemet{recherche augmentée}. Celui de l'accessibilité applicative et d'appuis à la production, pour une meilleure maîtrise de son environnement de travail, ce qu'aborde notre \gls{chatbot}. Et enfin celui de la production massive des données patrimoniales par un processus d'indexation semi-automatisée des \glspl{imgani}.\newline

Ces outils sont pensés pour \guillemet{révolutionner} ces trois axes par l'apport de solution simple, immédiate et fiable. Nous avons démontré que cela n'est pas réellement possible, ces outils étant faillibles et hautement évolutifs.  

Mais il faut surtout comprendre que de tels outils ne \guillemet{révolutionnent} par la gestion des collections patrimoniales mais impactent profondément diverses couches de cet univers métier. L'impact de ces outils n'est pas seulement technique, il est aussi juridique et éthique venant jusqu'à questionner et redéfinir des acquis.

Ainsi, pour chaque outil, nous allons aborder de front les questions soulevées par leurs implémentations dans les \gls{sgc}, au \coeur des pratiques et données des collections patrimoniales.

\subsection{La recherche augmentée}

Cet outil à une position intéressante car il s'agit, finalement, de transposer les pratiques de recherche grand public sur internet dans le logiciel. Ce qui dès lors crée des porosités entre notre position d'utilisateur d'internet, et celle d'utilisateur de logiciel métier. Nos pratiques publiques influent sur notre conception de concept que l'on retrouve dans le domaine professionnel, ce qui amène à un effort de création de lien entre les deux mondes afin d'unifier les pratiques.

Finalement cet outil amène à mélanger aux pratiques professionnelles des attentes \guillemet{grand public}. Aujourd'hui, lorsque l'on utilise un moteur de recherche, on formule instinctivement une phrase. Avec les \gls{ac}, c'est une dimension encore plus forte, qui entre de plus en plus dans nos habitudes. Ainsi, un mode de recherche traditionnel par mots-clés provoque un sentiment de frustration même au sein d'un logiciel métier. 

Cela pose donc un questionnement éthique important sur cette porosité entre l'usage public et l'usage professionnel. À quel moment l'infusion d'attentes et de pratiques \guillemet{grand public} dans des logiciels métier doit-elle s'arrêter ? On peut peux penser au moment où de telles pratiques mettent en danger les valeurs fondamentales de l'univers professionnel.\newline

Juridiquement, cela pose les entreprises fournisseuses de \gls{sgc}, comme SkinSoft, dans une position délicate. Elle propose un outil d'\gls{IA} traitant, sur son propre serveur, des données patrimoniales d'une institution tierce. Juridiquement, cela pose l'entreprise comme responsable de la conservation des données et de leur sécurité. Cela demande un contrôle total du traitement exécuté sur les données par le processus. 

Cette double responsabilité est encadrée par le \gls{RGPD}.\footnote{Voir \fullref{juridique}} Mais le développement de tels outils et leur popularité croissante dans ce contexte, peut amener à questionner sa pertinence. Ou plutôt son application à de tels processus demandant des redéfinitions juridiques pour un meilleur encadrement dans ce domaine. 

\subsection{Le \gls{chatbot}}

Pour le \gls{chatbot}, nous pouvons retrouver les mêmes considérations juridiques dès lors que nous le pensons comme assistant de contrôle de la qualité des données. Son accès aux données patrimoniales et sa capacité à les traiter posent les mêmes questions juridiques. 

Mais ce n'est pas l'aspect le plus important de cet outil. Les implications et enjeux sont bien plus d'ordre éthique et des pratiques métiers, que de l'ordre juridique.

Le premier élément est la redéfinition du rapport de l'utilisateur avec ses propres données. En effet, disposer d'un moyen de contrôle rapide, accessible et efficace peut-il impacter la productivité ? La possession d'un moyen de contrôle alimenté par de l'\gls{IA} peut pousser à croire à une infaillibilité de ce dernier et donc à baisser sa vigilance. Sans être une conséquence inévitable, c'est une considération éthique qu'il convient d'observer. 

Cela questionne directement le type d'outils que l'on doit et peut offrir dans un \gls{sgc}. Comment les concevoir afin de constituer un outil d'appui sans remplacer l'utilisateur humain ou créer un relâchement de sa vigilance ? Où se situe la frontière entre l'outil d'appui, et l'outil de remplacement ? 

Le \gls{chatbot} entre directement dans ces questions. Automatiser des tâches de contrôle et d'assistance par \gls{IA} ne risque-t-il pas de modifier trop profondément le rapport de l'utilisateur à son logiciel métier ? Trop de simplicité apportée par des outils provenant du quotidien peut-elle représenter un danger pour la productivité et la posture professionnelle des utilisateurs ? 

Nous ne saurons apporter de réponses. Cependant, il est certain que des réflexions éthiques et méthodologiques seront à mener vis-à-vis de ces outils. Il faut peut-être envisager d'encadrer les pratiques utilisateurs, en commençant par limiter l'impact de ces outils sur un équilibre professionnel préexistant.

\subsection{L'indexation semi-automatisée des \glspl{imgani}}

Mise à part des enjeux techniques et des questions d'\gls{ux} et d'\gls{ui}, cet outil remet en question des principes fondamentaux.\newline

D'abord le principe de rigueur scientifique. Contrairement à la \guillemet{recherche augmentée} ou au \gls{chatbot}, ici on génère de la donnée à valeur scientifique. Ce processus s’appuie sur des formes d'\gls{IAG} : les \gls{LLM}, \gls{VLM} et \gls{ALM}. La génération des données par ces outils peut être compromise par des hallucinations impactant la qualité des données et donc la rigueur scientifique.

Nous en revenons à la même problématique rencontrée avec la \guillemet{recherche augmentée}. La particularité ici est qu'une architecture \gls{RAG} n'est pas applicable. Cela contraint à adopter des architectures \textit{human in the loop} assez lourdes, pouvant compromettre le gain de temps que l'on cherche à obtenir. Il faut donc penser et établir de nouveaux procédés de contrôle des hallucinations afin de certifier une rigueur scientifique suffisante.

Ensuite, le principe de confiance. Fragiliser la véracité des informations des bases de données institutionnelles amène à provoquer des doutes dans l'usage de ces applications et dans les données produites. Ce qui, à plus large échelle, peut conduire à une méfiance du public auquel on diffuse ces informations.

On entre ici dans un domaine flou, à la croisée entre le procédé technique et les questions d'\gls{ux} et d'\gls{ui}. Habituellement, ces deux domaines permettent de rassurer l'utilisateur dans son usage d'une fonctionnalité. Mais devant la complexité de ces processus et les risques d'hallucinations, cela ne semble plus être suffisant. Il y a donc une demande de nouveaux moyens.\newline

Il se pose également la question de la provenance des données indexées. Les institutions conservatrices ont l'obligation de retracer la provenance des données conservées. On entre dans une limite juridique et éthique au sujet de l'\gls{IA}. Doit-on préciser que cela à été extrait par \gls{IA} ? Ou doit-on pousser l'\gls{IA} à citer les sources ? Le cadre juridique impose de préciser quand un contenu est généré par \gls{IA}\footnote{Voir \fullref{juridique}}. Mais lorsqu'il s'agit d'une donnée produite dans de tels processus, où l'\gls{IA} ne sert que d'extracteur, devons-nous le préciser ?  

Ainsi, il y a une redéfinition profonde des pratiques métiers dans la conservation du patrimoine. L'utilisateur du \gls{sgc} n'est plus créateur de la donnée, mais vérificateur et validateur ce qui change le paradigme.\newline

\subsection{Conclusion}

Finalement, l'implémentation de tels outils possède trois niveaux de complexités. Technologiques d'abord, car les problématiques, défis et attentes sont hétérogènes et complexes. De l'ordre éthique ensuite car leur implémentation et usage bouleversent des pratiques et valeurs portées par le monde de la gestion des collections. L’adaptation à ces éléments et nécessaire, mais il existe un fossé important pouvant demander une adaptation inversée. Et enfin juridique car ces outils peuvent être assez opaques et représenter des menaces face au principe de gouvernance des données. 

Nous sommes aujourd'hui au début de ces transformations profondes. Les outils ne sont pas encore mûr, et leurs utilisations professionnelles dans la gestion des collections pas encore effective. Cependant, elles annoncent d'ores et déjà des transformations importantes de cet univers professionnelles qu'il nous faut maîtriser et diriger.